% concatenated from main.tex
\documentclass[letterpaper, 10 pt, conference]{ieeeconf}  % Comment this line out if you need a4paper
\IEEEoverridecommandlockouts                              % This command is only needed if 
                                                          % you want to use the \thanks command

\overrideIEEEmargins                                      % Needed to meet printer requirements.


% See the \addtolength command later in the file to balance the column lengths
% on the last page of the document
\let\proof\relax 
  \let\endproof\relax
  \let\labelindent\relax
\usepackage{amsmath} % assumes amsmath package installed
\usepackage{amssymb}  % assumes amsmath package installed

\hyphenation{op-tical net-works semi-conduc-tor}
%\usepackage{showframe}  %Shows the page margins
\usepackage{tabulary}
\usepackage[english]{babel}
\usepackage{blindtext}
\usepackage[font=footnotesize]{caption}
\usepackage[makeroom]{cancel}
\usepackage{amsfonts}
\usepackage{pgfplots}
\usepackage{amsthm}  %This is for proofs
\usepackage{amssymb}
\usepackage{graphicx}
\usepackage{subcaption} % subfigures
\usepackage{pgf,tikz}
\usepackage{commath}
\pgfplotsset{ % Here we specify options for all figures in the document
  compat=newest, % Which version of pgfplots do we want to use?
   legend style =
  {font=\footnotesize },
  label style = {font=\footnotesize},
every tick label/.append style={font=\footnotesize}
  }
\usepackage{tabularx}
\usepackage{float}
\usepackage{cite}
\pagestyle{empty}
\usepgfplotslibrary{fillbetween}
  \usepgfplotslibrary{ternary}
\usetikzlibrary{shapes.multipart}
\usepackage{enumitem} % To reference numbered items
\usepackage{bbm,xurl} % This is for \mathbb for lower case letters


\usepackage[strict]{changepage}

\newcommand{\vect}[1]{\boldsymbol{#1}}
\newcommand{\mat}[1]{\boldsymbol{#1}}
\newcommand{\vt}[1]{\boldsymbol{#1}}
\newcommand{\wt}[1]{\widetilde{#1}}
\newcommand{\wh}[1]{\widehat{#1}}
\allowdisplaybreaks
\DeclareMathOperator{\sgn}{\text{sgn}}
\DeclareMathOperator{\sig}{\text{sig}}
\DeclareMathOperator{\diag}{\text{diag}}
\DeclareMathOperator{\spec}{\text{spec}}
\DeclareMathOperator*{\argmax}{\rm{argmax}} % thin space, limits underneath in displays

\newtheorem{remark}{Remark}
\newtheorem{theorem}{Theorem}
\newtheorem{lemma}{Lemma}
\newtheorem{assumption}{Assumption}
\newtheorem{observation}{Observation}
\newtheorem{definition}{Definition}
\newtheorem{corollary}{Corollary}
\newtheorem{proposition}{Proposition}
\newtheorem{property}{Property}
\newtheorem{problem}{Problem}
\newtheorem{example}{Example}

\renewcommand{\eqref}[1]{Eq.~(\ref{#1})}  %Modified equation reference

%\setlength{\abovecaptionskip}{0pt}
%\setlength{\belowcaptionskip}{-10pt}

\def\l {3.15}
\def\ll {6.5}
\def\qq {2.6}
\def\q {2.3}
\def\lll {2}
\def\h {1.9}
\def\hhh {1.6}
\def\hh {3.2}

\title{\LARGE \bf
Modelling the coevolution of opinion dynamics and decision making in social dilemmas}

% Power in Numbers: Modelling the Interdependence of Opinion Dynamics and Decision Making in Social Dilemmas.

% Modelling the coevolution of opinion dynamincs and decsion making in social dilemmas.

%

\author{Ella C. Davidson, Lorenzo Zino, Ming Cao, and Mengbin Ye
% <-this % stops a space
\thanks{E. C. Davidson and M. Ye are with the Adelaide Data Science Centre, Adelaide University, Adelaide, Australia (\texttt{\{connor.davidson;ben.ye\}@adelaide.edu.au}). L. Zino is with the Department of Electronics and Telecommunications, Politecnico di Torino, Torino, Italy (\texttt{lorenzo.zino@polito.it}). M. Cao is with the Engineering and Technology Institute Groningen, University of Groningen, Groningen, The Netherlands (\texttt{m.cao@rug.nl}). E. C. Davidson is supported by the Australian Government with a Research Training Program scholarship. M. Ye is supported by the Australian Government through the Australian Research Council (DE250100199) and Office of National Intelligence (NI240100203).}
}

\begin{document}

\maketitle
\thispagestyle{empty}

\begin{abstract}
This paper proposes a mathematical model for the coevolution of actions and opinions for a population facing a social dilemma. In particular, we assume each person participates in a Public Goods Game (PGG), with their action being to cooperate or defect, and holds an opinion about which action they prefer. We propose a payoff function that combines the PGG with the Friedkin--Johnsen model from opinion dynamics to form a coevolutionary game. According to a discrete-time process, players asynchronously update their actions and opinions, aiming to maximise their individual payoff for the coevolutionary game using myopic best-response. We study the equilibria and provide conditions for the existence of the all-defection and all-cooperation consensus equilibria. We also establish conditions for global convergence to the all-defection equilibrium.
%and numerical simulations identify a rich range of emergent behaviour of interest for the model, including persistence of cooperation and clustered scenarios. 
\end{abstract}

\section{Introduction}\label{sec:intro}
The emergence and persistence of widespread cooperation within a population when faced with a social dilemma is a fundamental problem of interest in biological, social and behavioural sciences~\cite{kollock_dilemmas_1998, dawes1980social}. This is partly because classical insights from game theory suggest that defection (also known as free riding) is often the best strategy from the perspective of a selfish rational individual. And yet, cooperation, which provides greater benefits for the entire population, is observed to persist in many real world scenarios. 

Collective action, i.e. the mobilisation of a large group to participate in a movement in pursuit of institutional or societal change, is a classical social dilemma~\cite{cantoni_protests_2024}. Participating in collective action is costly at an individual level and requires enough individuals to have a chance of success, but when successful, collective action often brings benefits to the entire population. More interestingly, there are cases of progressive movements (e.g. same-sex marriage and Black Lives Matter), in which many people participate even though they are not part of the repressed minority that would most benefit from the movements' success. 
% Widespread participation in these social movements (i.e. cooperation) yields benefits to society, but at the individual level, it is often easier not to participate (i.e. defect). %Another social dilemma of increasing relevance is democratic backsliding~\cite{carothers_dem_backslide_2024}. That is, a ``discontinuous series of incremental actions'' that erode a country's democratic processes, including the public's trust in democratic institutions~\cite{waldner_dem_backslide_2018}. Widespread participation in democratic activities such as voting (i.e. cooperation) yields benefits to the entire population, but at the individual-level, it is often easier to not participate (i.e. defect). 
% In many countries, there is an increasing level of defection, facilitated by individuals who do not believe in democratic processes~\cite{waldner_dem_backslide_2018}. 
% While many notable countries are experiencing this phenomenon, like Hungary, Tunisia, and the United States, the power to overturn such effects lies with the people \cite{carothers_dem_backslide_2024}. Attacks on Poland’s democracy sparked large-scale anti-government protests, which led to an alternation of power \cite{carothers_dem_backslide_2024}. Such a scenario prompts the question, what drives individuals to participate in a social movement?

% Collective action, i.e. the mobilisation of broad parts of a population to participate in a movement in pursuit of institutional or societal change, is a classical social dilemma. Participating in protests and social movements can be costly at an individual level \cite{cantoni_protests_2024}. However, if enough individuals protest/cooperate, the movement can succeed~\cite{biggs_protest_size_2018}, benefiting the entire population. Of particular note are movements that led to progressive change (e.g. the legalisation of same-sex marriage, women's suffrage, or Black Lives Matter), many people who joined the movement were not part of the repressed minority that would benefit from the change. 

% In contrast to collective action and protest, another social dilemma of increasing concern in many countries is democratic backsliding~\cite{carothers_dem_backslide_2024}. Democratic backsliding is generally defined as a ``discontinuous series of incremental actions'' that erode a country's democratic processes, including the public's trust in democratic institutions~\cite{waldner_dem_backslide_2018}. This dilemma is facilitated and accelerated by individuals deciding not to participate in democratic activities like voting, essentially opting to defect from democratic participation~\cite{waldner_dem_backslide_2018}. While many notable countries are experiencing this phenomenon, like Hungary, Tunisia, and the United States, the power to overturn such effects lies with the people \cite{carothers_dem_backslide_2024}. Attacks on Poland’s democracy sparked large-scale anti-government protests, which led to an alternation of power \cite{carothers_dem_backslide_2024}. Such a scenario prompts the question, what drives individuals to participate in a social movement?

%Various factors that allow cooperation to emerge and persist over time have been proposed and studied. Peer punishment is often given to, or threatened to, individuals who deviate from a group's norm~\cite{abbink_punish_2017}. Reciprocity, where cooperation is rewarded with friendliness and further cooperation, can further facilitate cooperation~\cite{gouldner_reciprocity_1960, park_opinions_2021}. Social influence from other people can be a significant role in how an individual acts, and thus the structure of a social network can facilitate cooperation if a cluster of individuals interact more often with each other compared to individuals outside of the cluster~\cite{nowak_cooperation_2006}.

Various factors that allow cooperation to emerge and persist over time have been proposed and studied. These include, and are not limited to, peer punishment~\cite{abbink_punish_2017}, reciprocity~\cite{gouldner_reciprocity_1960, park_opinions_2021}, and social influence~\cite{nowak_cooperation_2006}. While these factors are social in nature, a key psychological factor in determining a person's action is their own belief~\cite{henry2003beyond}. This is true for many contexts, including especially social dilemmas and collective action, as observed empirically within the literature on the social psychology of protest movements~\cite{thomas_protest_2024}. Moreover, generally, a protester's beliefs, attitudes, and psychological drivers (all of which evolve as the movement gains or loses traction) are central to their engagement with a movement~\cite{louis2022failure,nardini_BLM_2021}.

% We posit that individuals participate in a social movement, or cooperate when faced with a social dilemma, because their actions and beliefs are intrinsically connected—an individual chooses to contribute because they believe it is the right thing to do. 
 
% Recent research has also highlighted the importance of a shared social identity among those participating in a movement~\cite{thomas2020testing}. 

Despite empirical evidence of an interdependence between an individual's action and opinion, there have been limited efforts to combine the two into a unified mathematical model. The continuous-opinion discrete-action (CODA) model framework provided important initial contributions~\cite{martins_coda_2008,ceragioli_coda_2018}, operating under the assumption that an individual's action is a direct quantisation of their opinion, limiting the possibility of capturing the fundamental tension between cooperation and defection in social dilemmas. More recently, a mathematical model was proposed using a unified game-theoretic formulation, in which individuals played coordination games on a network and exchanged opinions on preferred strategies~\cite{aghbolagh_coevolutionary_2023}. The resulting coevolution of actions and opinions was studied analytically, and emergent phenomena such as unpopular norms and polarisation were identified~\cite{aghbolagh_coevolutionary_2023}. 
The model was extended by~\cite{liang_coevolution_2025} to consider an anti-coordination game instead. 
% for agents playing a network coordination game on a two-layer network. One layer allowed agents to exchange opinions and influence one another to change the outcome of the coordination game played on the other layer of the network~\cite{aghbolagh_coevolutionary_2023}. 
%The model was extended by~\cite{liang_coevolution_2025} to allow individuals to partake in an anti-coordination game instead. 
% However, the models above have yet to be utilised to study cooperation in social dilemmas.

% When it comes to modelling the context of collective action, a public goods game is an appropriate choice. The PGG posits that $n$ players can decide to either contribute to the public good (to cooperate) or not (to defect)~\cite{govaert_pgg_2021}. This general concept can be formulated in many ways, from each agent deciding to contribute a certain amount from a personal store or deciding whether to simply cooperate/defect~\cite{govaert_pgg_2021}. The latter definition makes more sense within the context of collective action, as the decision to participate in a protest can be costly at an individual level~\cite{cantoni_protests_2024}. Previously, a model proposed by~\cite{park_opinions_2021} paired a PGG with opinion dynamics. In the model, agents decided how much to contribute to the public good and the agents’ opinions were formed from the PGG payoff~\cite{park_opinions_2021}. The modelling of~\cite{park_opinions_2021} illustrates a one-way relationship of actions affecting opinions, and not a coevolution of the intertwined relationship of the two, signalling the need for a more intricate model to be created.

In this paper, we consider the coevolutionary dynamics of actions and opinions of a group of people in the context of social dilemmas. 
% of players participating in a Public Goods Game (PGG), a classical game used in game theory to study social dilemmas due to its flexibility and broad applicability~\cite{Bramoull2007}. 
In more detail, individuals participate in a Public Goods Game (PGG), a classical game commonly used for modelling social dilemmas due to its flexibility and broad applicability~\cite{Bramoull2007}, and decide whether to cooperate or defect. An individual's strategic decision to cooperate or defect coevolves with their opinion (or attitude/belief) towards cooperation. In \cite{park_opinions_2021}, a PGG was integrated within a continuous-time opinion dynamics model, yielding a framework similar to CODA models. Players revise their opinion based on a realistic opinion dynamics model that accounts for social interactions and the payoff of a PGG, and then make decisions on their actions (the amount to contribute to the public good) based on their opinion in a stochastic manner. This, however, illustrates a one-way relationship, not an interdependent coevolution of actions and opinions.

%In \cite{park_opinions_2021}, a PGG was integrated within a continuous-time opinion dynamics model, building on a framework similar to CODA models. Players revise their opinion based on a realistic opinion dynamics model that accounts for social interactions and the payoff of a PGG, and then make decisions on their actions (the amount to contribute to the PGG) based on their opinion in a stochastic manner. This, however, illustrates a one-way relationship, limiting the possibility to capture misalignments between actions and opinions that are often observed in real life.


% The proposed model effectively unites two streams of literature---opinion dynamics and evolutionary game theory.
Here, we adopt the game-theoretic coevolutionary approach~\cite{aghbolagh_coevolutionary_2023} and propose a joint payoff function for the coevolutionary game that combines i) the payoff from the PGG, ii) opinion dynamics evolving through social influence (inspired by the Friedkin--Johnsen model~\cite{friedkin_social_1990}), and iii) a term that captures an individual's desire for consistency between their action and opinion. We consider a population of players interacting on a network to share opinions, revising their actions and opinions asynchronously at discrete timesteps using myopic best-response updating. Our work contrasts the coevolutionary models in \cite{aghbolagh_coevolutionary_2023,liang_coevolution_2025}, which consider coordination and anti-coordination games rather than social dilemmas, and offers a novel lens to study the emergence and persistence of cooperation in social groups.

The main contributions of the paper are as follows. We propose a novel coevolutionary game, and obtain the explicit expression for the best-response updating, showing how a player's best-responding action and opinion are shaped by others' opinions and the PGG parameters. We study the equilibria, providing sufficient conditions for the existence of the all-defection and all-cooperation equilibria, and identify a sufficient condition for the all-defection equilibrium to be the unique equilibrium with a global stability property. The rest of the paper is organised as follows. Section~\ref{sec:model} presents the coevolutionary game and associated dynamics. Section~\ref{sec:results} contains the main results and Section~\ref{sec:conclusion} concludes the paper.
%Numerical simulations illustrate the theoretical results, as well as showcasing equilibria with a mixture of defectors and cooperators.
% .  i) computation of a player’s best-response given the intertwined mechanisms; ii) analysis of equilibria of interest for the model; and iii) a rigorous convergence result for a special case of the dynamics. We also include a section exploring numerical simulations of the model, showcasing equilibria not foreseen through our mathematical analysis. 

 %with numerical simulations in Section~\ref{sec:simulations}.
% , we explore numerical simulations of the coevolutionary model. 


\section{The Model}\label{sec:model}

\subsection{Preliminaries}
% \subsubsection*{Notation}
The set of real and non-negative integer numbers are denoted by $\mathbb{R}$ and $\mathbb{Z}_+$ respectively. Bold lower case and upper case fonts denote a vector $\vect{x}$ and matrix $\vect{A}$, respectively, with the $i^{th}$ component being $x_i$ and the $j^{th}$ entry in the $i^{th}$ row being $a_{ij}$, respectively. The column vector of all ones and all zeros is denoted by $\vect{1}$ and $\vect 0$, respectively, and the identity matrix is $\vect{I}$; the appropriate dimension is understood within the context of its use. 
% Given a vector $\vect{x}$ or matrix $\vect{A}$, the transpose is defined as $\vect{x}^\top$ and $\vect{A}^\top$ respectively.

A weighted network is defined as a tuple $\mathcal{G} = (\mathcal{V}, \mathcal{E}, \mat W)$, where $\mathcal{V}=\{1,\dots,n\}$ is the set of $n$ nodes; $\mathcal{E}\subseteq\mathcal{V}\times\mathcal{V}$ is the edge set, whereby $(i,j)\in\mathcal{E}$ iff there is a link going from node $i$ to $j$, and $\boldsymbol{W}\in\mathbb{R}_+^{n\times n}$ is the weight matrix, so that entry $w_{ji}>0\Longleftrightarrow (i,j) \in \mathcal{E}$. A network is undirected iff $\vect{W}=\vect{W}^\top$. We assume that $\mat W$ is \textit{row-stochastic}, i.e., the entries of each row add to $1$, that is $\vect{W1}=\vect 1$.

A game $\Gamma(\mathcal{V}, \mathcal{A}, \boldsymbol{\pi})$ is defined by a set of players $\mathcal{V}=\{1,\dots,n\}$; a set of strategies $\mathcal{A}$ that the players can choose (here assumed to be the same for all agents); and a payoff vector $\boldsymbol{\pi}=[\pi_i, \dots, \pi_n]^\top$, where $\pi_i:\mathcal{A}^n\rightarrow\mathbb{R}$ is the payoff function for player~$i\in \mathcal{V}$. We let $z_i \in \mathcal{A}$ be the strategy of player $i$, which can be a scalar or vector and discrete or continuous. The vector $\vect{z}=[z_1,\dots,z_n]\in\mathcal A^{n}$ gathers the strategies of all players. Given a player $i\in\mathcal V$, we denote by $\vect{z}_{-i}=[z_1,z_2,\dots,z_{i-1},z_{i+1},\dots,z_n]\in\mathcal A^{n-1}$ the vector of the strategies of all players, except for player~$i$. The payoff that player~$i$ receives for selecting a strategy $z_i\in\mathcal A$ is typically dependent on the action of the others $\vect{z}_{-i}$, so it is denoted as $\pi_i(z_i,\vect{z}_{-i})$. 
Given a game, we define the set of best-response strategies and a Nash equilibrium as follows.
\begin{definition}[Best response and Nash Equilibrium]\label{def:BR}
    Given a game $\Gamma(\mathcal{V}, \mathcal{A}, \boldsymbol{\pi})$, the best-response strategies for player~$i\in \mathcal{V}$, given the state of the system $\vect{z}\in\mathcal A^n$ are defined as 
    \begin{equation}\label{eq:best_response}
        \mathcal{B}_i(\pi_i(\cdot,\vect{z}_{-i})):= \argmax\nolimits_{{z_i}\in \mathcal{A}} \ \vect \pi (z_i,\vect{z}_{-i}).
    \end{equation}
    A Nash equilibrium $\vect{z}^*=[z_1^*,\dots,z_n^*]\in\mathcal A^{n}$ is a strategy profile such that $z_{i}^* \in \mathcal{B}_i( \pi_i(\cdot,\boldsymbol{z}^*_{-i})),\, \forall\,i\in\mathcal V$.
\end{definition}

% \begin{definition}[Nash equilibrium]
%     Given a game $\Gamma(\mathcal{V}, \mathcal{A}, \boldsymbol{\pi})$, a Nash equilibrium $\vect{z}^*=[z_1^*,\dots,z_n^*]\in\mathcal A^{n}$ is a strategy profile such that $z_{i}^* \in \mathcal{B}_i( \pi_i(\cdot,\boldsymbol{z}^*_{-i})),\, \forall\,i\in\mathcal V$.
%     % \begin{equation}\label{eq:nash}
%     %     z_{i}^* \in \mathcal{B}_i( \pi_i(\cdot,\boldsymbol{z}^*_{-i})),\quad \forall\,i\in\mathcal V.
%     % \end{equation}
% \end{definition}

\subsection{The PGG and Actions}

% \begin{figure}
%     \centering
%     \includegraphics[width=1\linewidth]{fig/public_goods_game.pdf}
%     \caption{Schematic of a simple public goods game.}
%     \label{fig:pgg}
% \end{figure}

We consider the set of players $\mathcal V$ participating in a PGG. Each player can decide whether to contribute to the public good or not, referred to as ``cooperate" or ``defect", respectively. We denote player $i$'s decision with the term $x_i \in \{0,1\}$, with $x_i=0$ and $x_i=1$ representing player~$i$ deciding to defect and cooperate, respectively. All $n$ player decisions are held in the action vector $\vect{x}=[x_1, \dots, x_n]^\top \in \{0,1\}^n$. 

If player~$i$ decides to cooperate, they contribute a unit amount to the public good. Regardless of whether they have cooperated or not, each player receives a reward, with the amount depending on how many players have cooperated. 
% Specifically, let $\xi=\frac{1}{n}\sum_{j=1}^nx_j$ be the fraction of players who decide to cooperate. Then, we define a function $\rho(\cdot):[0,1]\to\mathbb R^+$ such that the total reward of the PGG is $\rho(\xi)$, which is split evenly across the players (i.e., each player receives exactly $\frac1n\rho(\xi)$). In the most common scenario, we can assume that such a function is linear, i.e., we let $\rho(\xi)=r\xi$ where the constant $r$, with $1<r<n$, represents the public good multiplier. This variable can also be seen as the benefit-to-cost ratio of the game. 
Thus, a player's payoff depends on the player's own action and the actions of the other $n-1$ players. Here, we consider a simple implementation of the PGG, with all players contributing a unit amount to the same public good if they cooperate, and a constant $r$, with $1<r<n$, that represents the public good multiplier~\cite{Bramoull2007, govaert_rationality_2021}. Thus,
% Hence, player~$i$'s payoff for the PGG is equal to 
\begin{align}\label{eq:payoff_PGG}
    \pi_i^{a}(x_i,\vect{x}_{-i}) &= x_i \Bigg[  \dfrac{r\left(\sum^n_{j=1, j\neq i}x_j + 1\right)}{n} -1 \Bigg] \notag \\ 
    &+ (1-x_i)\dfrac{r}{n}\sum\nolimits^n_{j=1, j\neq i}x_j,
\end{align}
where the first term (which is non-zero if player~$i$ cooperates) accounts for the share of the total reward received by player~$i$ if they cooperate, less their unit contribution), and the second term (which is non-zero if player~$i$ defects) accounts for the share of the total reward received by player~$i$ if they defect. 
% A simple schematic of a PGG is shown in Fig.~\ref{fig:pgg}.

% \begin{remark}
% Potential extensions of this framework can consider 1) non-linear reward functions such as a threshold function or 2) network public goods games, whereby agents do not play in an all-to-all setting. 
% \end{remark}

\subsection{Opinion Dynamics}
To model the evolution of the players' opinions we opt for the use of the empirically validated Friedkin--Johnsen (FJ) model~\cite{friedkin_social_1990}. Each player~$i$ of a population $\mathcal V=\{1,\dots,n\}$ is assigned a continuous random variable $y_i\in [0,1]$, which represents player~$i$'s opinion. All $n$ of the players' opinions are held in the opinion vector $\vect{y}=[y_1, \dots, y_n]^\top \in [0,1]^n$. Players share their opinion on a social network $\mathcal G=(\mathcal V, \mathcal E,\mat W)$, where the entry $w_{ij}$ of $\mat W$ represents the social influence of player~$j$ on player~$i$. As detailed in the literature~\cite{proskurnikov_tutorial_2017}, the FJ model can be formulated as a game~\cite{ghaderi_opinion_2014}, where player~$i$'s payoff for adopting an opinion ${y}_i$, given an opinion vector $\vect y$ is defined as
\begin{equation}\label{eq:opinion}
    \pi_i^{o}(y_i,\vect{y}_{-i}) = -\dfrac{1}{2}(1-\gamma_i)\sum_{j \in \mathcal{V}}w_{ij}(y_i - y_j)^2 - \dfrac{1}{2}\gamma_i(y_i - u_i)^2,
\end{equation}
where the parameter $\gamma_i\in[0,1]$ represents player~$i$'s level of attachment to their constant prejudice $u_i\in[0,1]$. The best-response to this payoff yields the FJ opinion dynamics model~\cite{friedkin_social_1990}. Note that the FJ model, as standard in opinion dynamics, allows $w_{ii} > 0$ (self-loops in $\mathcal G$), but in a game-theoretic context corresponds to the non-standard concept of playing a game against oneself. This can be addressed by a minor adjustment to \eqref{eq:opinion}, see \cite{aghbolagh_coevolutionary_2023}. 
% In this context, self-loops add social inertia to the dynamics. %In this case, the summand involving $w_{ii}$ should be written as $w_{ii}(y_i-\tilde y_i)^2$, where we treat $\tilde y_i$ as the opinion of another (virtual) player.

\subsection{Coevolutionary Model}
%Our coevolutionary model consists of a coevolutionary game, and a dynamics in which players update their strategies by performing a best-response. In the following, we describe each component separately.

\subsubsection*{Coevolutionary Game} We are now in a position to introduce the coevolutionary game that combines the public goods and opinion dynamics game in an interdependent manner. In this setting, each player plays a PGG, where they can decide to defect, $x_i=0$, or to cooperate, $x_i=1$. Simultaneously, each player holds an opinion, $y_i \in [0,1]$, representing their level of support (or preference) for the action of cooperating in the PGG. Specifically, a player~$i$ fully supporting (preferring) defection and cooperation correspond to $y_i=0$ and $y_i=1$, respectively. We make the important distinction between whether a player chooses to defect or cooperate ($x_i$) and whether a player supports it ($y_i$), noting that these two may not always be aligned. Player~$i$'s strategy is defined as $z_i = (x_i, y_i)$, and we refer to the vector $\vect{z}= (\vect x, \vect y)$ as the \textit{state of the system} and strategy profile interchangeably.

% We now introduce the payoff function for the coevolutionary game. 
Given the state of the system $\boldsymbol{z}=(\boldsymbol{x},\boldsymbol{y})$, we propose that the coevolutionary game has the following payoff function for player~$i$ selecting strategy $z_i = (x_i, y_i)$:
\begin{align}\label{eq:full_payoff}
    \pi_i(z_i,\vect{z}_{-i}) &= \alpha_i\pi_i^{a}(x_i,\vect{x_{-i}}) + \beta_i\pi_i^{o}(y_i,\vect{y_{-i}}) \notag \\ 
    & \qquad - \dfrac{1}{2}\lambda_i(x_i-y_i)^2,
\end{align}
with $\pi_i^{a}(x_i,\vect{x}_{-i})$ and $\pi_i^{o}(y_i,\vect{y}_{-i})$ from \eqref{eq:payoff_PGG} and \eqref{eq:opinion}, respectively. 
The parameters $\alpha_i,\beta_i,\lambda_i \in [0,1]$ represent player~$i$'s relative weight of actions, opinions, and desire to ensure consistency between their action and opinion (self-consistency), respectively. Without any loss in generality, we assume that $\lambda_i=1-\alpha_i-\beta_i$, from which \eqref{eq:full_payoff} can be readily interpreted as a convex combination of the three terms, one term proportional to the payoff of the corresponding PGG, one term proportional to the payoff of the corresponding opinion dynamics, and a third term that accounts for self-consistency. A major conceptual foundation of our model is to explicitly distinguish between action and opinion, while introducing the self-consistency term to reduce the total payoff as the player's action and opinion become less aligned.  All player payoffs for the coevolutionary game are stored in the vector $\vect{\pi}$, i.e. $\vect{\pi} = [\pi_1, \cdots, \pi_n]^\top$. We can now formally define the coevolutionary game as follows.

\begin{definition}[Coevolutionary game] \label{d:coev_game}
    The coevolutionary game is $\Gamma(\mathcal{V}, \mathcal{A}, \boldsymbol{\pi})$, played by the set of players $\mathcal V$, on $\mathcal G = (\mathcal{V},\mathcal E, \mat W)$, with strategy set $\mathcal A = \{0,1\}\times[0,1]$ and payoff vector function $\vect{\pi}$, where entry $\pi_i$ is defined in~\eqref{eq:full_payoff}.
\end{definition}

In order to characterise the best-response strategies to the coevolutionary game, we introduce a \textit{discriminant} quantity:%$\delta_i$ as a function of the state $\vect z$:
\begin{align}\label{eq:discriminant}
    \delta_i(z_i,\vect{z}_{-i}) &=  \alpha_i\left( \dfrac{r}{n}-1 \right) \\ \notag
    &+ \frac{\beta_i\lambda_i}{\beta_i+\lambda_i}\bigg(\gamma_iu_i+(1-\gamma_i)\sum_{j\in\mathcal{V}}w_{ij}y_j-\frac{1}{2}\bigg).
\end{align} 

The following result provides the best-response explicitly, by utilising the discriminant $\delta_i$. The proof is in the Appendix.
\begin{proposition}\label{prop:BR}
    Consider the coevolutionary game from Definition~\ref{d:coev_game}. Then,  %best-response strategy for player $i$, 
    given a state $\vect z\in\mathcal A^n$, for a player $i\in\mathcal V$, $(\bar x_i,\bar y_i)\in \mathcal{B}_i(\vect \pi_i(\cdot,\boldsymbol{z}_{-i}))$ if and only if
    \begin{align}
        \bar{x}_i &= \hat s_i(\delta_i(\vect{y})) \label{eq:BR_action} \\
        \bar y_i &= \dfrac{\beta_i(1-\gamma_i)\sum_{j\in\mathcal{V}}w_{ij}y_j+\beta_i\gamma_iu_i+\hat s_i(\delta_i(\vect y))\lambda_i}{\beta_i+\lambda_i},\label{eq:BR_opinion}
    \end{align}
    with %$\bar y_i \in [0,1]$ and
    \begin{equation}
        \hat s_i(\delta_i(\vect{y})) = \left\{\begin{array}{ll}0&\text{if }\delta_i(\vect{y})\leq 0\\
        1&\text{if }\delta_i(\vect{y})\geq0.\end{array}\right.
    \end{equation}    
\end{proposition}

\begin{remark}\label{rem:br}When $\delta_i = 0$, the set of best responses comprises two elements: one with $\bar x_i=0$ and one with $\bar x_i=1$, and the corresponding values of $\bar y_i$, which differ, as per \eqref{eq:BR_opinion}.%both $(1,y_i^{1})$ and $(0,y_i^{0})$ are best-responses
\end{remark}
%, and we make a mild assumption that the player in this instance will opt to defect.
Finally, while the $\delta_i$ is presented in \eqref{eq:discriminant} as a function of $\vect z = (\vect x, \vect y)$, it turns out it is in fact only a function of $\vect y$. This property will play a key role in our analysis of the updating dynamics in the sequel. We note that~\eqref{eq:discriminant} suggests a one-way relationship between the players' opinions and actions, as we observed in~\cite {park_opinions_2021}. However, the formulation of the payoff in~\eqref{eq:full_payoff} does not. We conjecture that this results from the linear payoff formulation of the PGG; if a non-linear PGG payoff were implemented, this may no longer be the case.
% {\color{magenta}
% \begin{remark}
%     We note that~\eqref{eq:discriminant} suggests a one-way relationship between the players' opinions and actions, as we observed in~\cite {park_opinions_2021}. However, the formulation of the payoff in~\eqref{eq:full_payoff} does not. We conjecture that this results from the linear payoff formulation of the PGG; if a non-linear PGG payoff were implemented, this may no longer be the case.
% \end{remark}
% }

\subsubsection*{Player Dynamics}
Let us now define the update dynamics for the players, who are repeatedly playing the coevolutionary game. At each discrete timestep $t \in \mathbb{Z}_+$, we posit that a set $\mathcal{R}(t) \subseteq \mathcal{V}$ of players becomes active and revises their strategy, and $\mathcal R(t)$ satisfies the following property.

\begin{assumption}\label{assumption:revision_sequence}
    (Revision sequence). There exists a constant $T<\infty$ such that $\cup_{s=0}^{T-1}\mathcal{R}(t+s) = \mathcal{V},$ for any $t \geq 0$.
\end{assumption}

\begin{remark}
    Assumption~\ref{assumption:revision_sequence} covers many synchronous and asynchronous update rules. Synchronous updating corresponds to $\mathcal{R}(t) = \mathcal{V}$ for all $t$, while asynchronous updating corresponds to $|\mathcal{R}(t)| = 1$, for all $t$. The condition of $\cup_{s=0}^{T-1}\mathcal{R}(t+s) = \mathcal{V}$ mimics typical assumptions in opinion dynamics with time-varying networks~\cite{proskurnikov_tutorial_2017}, and ensures that each player activates at least once within a time window of $T$, and that this occurs infinitely often as $t\to\infty$. 
\end{remark}

We assume that active players in $\mathcal R(t)$ will always update their future action and opinion by adopting the best-response strategy for the coevolutionary game given $\vect z(t)$. Hence, at timestep $t$, each player~$i \in \mathcal{R}(t)$ updates their state as
\begin{equation}
    z_i(t+1) \in \mathcal{B}_i(\pi_i(\cdot,\boldsymbol{z}_{-i}(t)))
\end{equation}
where $\mathcal{B}_i$ is the set of best-responses strategies computed in Proposition~\ref{prop:BR}. Here, we make a mild assumption that if a player's best response comprises two elements (see Remark~\ref{rem:br}), they will opt to defect due to inherit selfishness. Therefore, based on Proposition~\ref{prop:BR}, an active player~$i$ with state $z_i(t)=(x_i(t),y_i(t))$ is updated as 
\begin{align}
    x_i(t+1) &= s_i(\vect{y}(t)) \label{eq:update_action} \\
    y_i(t+1) &= \dfrac{\beta_i(1\!-\!\gamma_i)\sum_{j\in\mathcal V} w_{ij}y_j(t)\!+\!\beta_i\gamma_iu_i\!+\!s_i(\vect y(t))\lambda_i}{\beta_i+\lambda_i} \label{eq:update_opinion}
\end{align}
where
\begin{equation} \label{eq:update_s}
    s_i(\vect{y}(t)) = \left\{\begin{array}{ll}0&\text{if }\delta_i(\vect{y}(t))\leq 0\\
    1&\text{if }\delta_i(\vect{y}(t))>0,\end{array}\right.
\end{equation}
with $\delta_i(\vect{y}(t))$ defined in \eqref{eq:discriminant}. Any player $j \notin \mathcal{R}(t)$ is inactive, and thus $z_j(t+1)=z_j(t)$. 

We now point out some important features of the player update dynamics, and relevant interpretations within the context of the modelling problem. 

First, notice from \eqref{eq:update_action} that $x_i(t+1)$ is determined solely by $\delta_i$, which is a function of $\vect y(t)$. In other words, the players decision to defect or cooperate can potentially (but not necessarily) be influenced by other players' opinions, but not actions. However, $\delta_i$ also contains the parameters $r,n,\alpha_i$ which are associated with the PGG, and as we illustrate in the sequel, there are parameter values for which the sign of $\delta_i$ is independent of $\vect y$. This is a key departure from the coevolutionary model in \cite{aghbolagh_coevolutionary_2023}, which considered coordination games on networks with Friedkin--Johnsen opinion dynamics, and the equivalent $\delta_i$ that determined a player's action was dependent on both $\vect x$ and $\vect y$.

Second, a player's new opinion is a convex combination of their neighbours' current opinions, their existing prejudice, and their new action. This is a natural extension of the Friedkin--Johnsen model, which considers a convex combination of the former two quantities. It also illustrates how a player's new opinion is impacted by both their neighbours' current opinions and their own decision, illustrating the coevolving nature of the dynamics. Because $\alpha_i+\beta_i+\lambda_i = 1$, the dynamics are clearly influenced by the relative weighting that a player places on the PGG ($\alpha_i$), the opinion dynamics ($\beta_i$) and the need to be self-consistent ($\lambda_i$).

In summary, the \textit{coevolutionary game} $\Gamma=(\mathcal{V}, \mathcal{A}, \vect{\pi})$ is defined on a network $\mathcal{G}$ as in Definition~\ref{d:coev_game}. The \textit{coevolutionary model} is the dynamical system arising from the game being played repeatedly over discrete time-steps by the players, in which active players at each time-step update their action and opinion using a best-response approach. 

\section{Main Results}\label{sec:results}
Here, we study the equilibria of the system, characterizing them as Nash equilibria of the coevolutionary game, discussing the existence of some equilibria of interest (namely, those in which the entire population defect or cooperate, respectively), and show that under certain parameter conditions, there is global convergence to a unique equilibrium.

\subsection{Classification of Equilibria}
To characterise the equilibria of the coevolutionary model, we introduce some useful terminology. A state $\vect z^* = (\vect x^*, \vect y^*)$ is an equilibrium if \eqref{eq:update_action} and \eqref{eq:update_opinion} hold with $x_i(t+1) = x_i(t) = x_i^*$ and $y_i(t+1)=y_i(t)=y_i^*$ for all $i\in\mathcal V$.

Our results begin by relating the equilibria of the coevolutionary model to Nash equilibria of the coevolutionary game. We omit the proof for brevity, noting it can be obtained directly from Definition~\ref{d:coev_game} and Eqs.~(\ref{eq:update_action})--(\ref{eq:update_s}).

\begin{proposition}
   Under Assumption~\ref{assumption:revision_sequence}, the set of equilibria of the dynamics for the coevolutionary model in Eqs.~(\ref{eq:update_action})--(\ref{eq:update_s}) coincides with the set of Nash Equilibria of the coevolutionary game in Definition \ref{d:coev_game}. 
\end{proposition}

% {\color{magenta}\begin{proof}
%     This proof is omitted due to page restrictions.
% \end{proof}}

% \begin{proof}
%     If $\vect z^*\in\mathcal A^n$ is a Nash equilibrium, then by definition $\vect z^*\in\mathcal{B}(\vect z^*), \forall i \in \mathcal{V}$, which implies $\vect z(t+1)=\vect z(t) = \vect z^*, \forall t$. If $\vect z^*$  is an equilibrium of the coevolutionary model, then we have $\vect z(t)=\vect z(t+1)= \ldots =\vect z(t+T) = \vect z^*$ with $T$ defined in Assumption~\ref{assumption:revision_sequence}. This  implies no player has revised their strategy, i.e. they are all playing a best-response strategy. Thus, $ z_i^*\in\mathcal{B}_i(\pi(z_i^*,\vect z^*_{-i})), \forall i \in \mathcal{V}$.
% \end{proof}

Henceforth, our analysis is written from the perspective of characterising the equilibria for the coevolutionary model, but as indicated above, this provides equivalent characterisation of the Nash equilibria for the coevolutionary game. In this paper, we derive results under the following assumption.
\begin{assumption}\label{a:0}
    All players have zero prejudice, i.e. $\gamma_i = 0, \forall i \in \mathcal{V}$, and $\alpha_i, \beta_i, \lambda_i \in (0,1), \forall i \in \mathcal{V}$.
\end{assumption}
This assumption allows us to the focus on the role of the variables $\alpha_i, \beta_i,$ and $\lambda_i$, which are the parameters of the PGG, opinion dynamics, and self-consistency. Extending the results to allow for $\gamma_i > 0$ is an important future direction, to examine the impact of attachment to existing prejudice.
% {\color{magenta}
% \begin{remark}
%     We implement Assumption~\ref{a:0}, specifically $\gamma_i = 0$, for all $i \in \mathcal{V}$, for analytical convenience, as this simplifies the computations and clarifies the interplay between the variables $\alpha_i, \beta_i,$ and $\lambda_i$. Extensions incorporating $\gamma_i$ offer a potential avenue for future work to understand its influence better.
% \end{remark}
% }

% We now record two expressions that are essential to our analysis. Consider player $i$ who becomes active at timestep $t$, i.e., $i\in\mathcal R(t)$. For them to decide to defect at the next timestep, we must have $\delta_i(z_i(t),\vect{z}_{-i}(t)) \leq 0$, which we can prove using \eqref{eq:discriminant} that is equivalent to
% \begin{subequations}\label{eq:delta_cond}\begin{equation}\label{eq:delta_neg_cond}
%         \sum_{j\in\mathcal{V}}w_{ij}y_j(t) \leq \dfrac{-\alpha_i(1-\alpha_i)\left(\dfrac{r}{n}-1\right)}{\beta_i\lambda_i}+\frac{1}{2}.
% \end{equation}
% Conversely, for player $i$ to cooperate at the next timestep, then $\delta_i(z_i(t),\vect{z}_{-i}(t)) > 0$, which is equivalent to
% \begin{equation}\label{eq:delta_pos_cond}
%     \sum_{j\in\mathcal{V}}w_{ij}y_j(t) > \dfrac{-\alpha_i(1-\alpha_i)\left(\dfrac{r}{n}-1\right)}{\beta_i\lambda_i}+\frac{1}{2}.
% \end{equation}\end{subequations}

To better interpret the results, let us define three types of consensus that can arise. 

\begin{definition}[Classification of consensus states]
    A state $\vect z = (\vect x, \vect y)$ is said to be an {action consensus} if $x_i=x_j \ \forall i,j \in \mathcal{V}$. For any arbitrary $\vect y$, we call the action consensus state $\vect z = (\vect 0, \vect y)$ as the all-defection state, and $\vect z = (\vect 1, \vect y)$ as the all-cooperation state. Similarly, we define an {opinion consensus} as $y_i=y_j \ \forall i,j \in \mathcal{V}$. Then, for an arbitrary $\vect x$, the state $\vect z = (\vect x, c\vect 1)$ is an opinion consensus state for any $c\in [0,1]$. Finally, we call $\vect z = (\vect 0, \vect 0)$ the {all-defection consensus} state, and $\vect z = (\vect 1, \vect 1)$ the {all-cooperation consensus} state.
\end{definition}

We begin our equilibria analysis by showing that any equilibrium in which there is action consensus must in fact have an opinion consensus.

\begin{proposition}\label{prop:consensus}
    Consider the coevolutionary model under Assumption~\ref{a:0}, and let $\vect{z}^*=(x^*,y^*)$ be an equilibrium of the model. There holds $\vect x^* = \vect 0 \Rightarrow \vect y^* = \vect 0$ and $\vect x^* = \vect 1 \Rightarrow \vect y^* = \vect 1$. That is, an equilibrium $\vect z^*$ with an action consensus must in fact also have an opinion consensus.  
\end{proposition}
\begin{proof}
    We first show that $\vect{x}^*= \vect{0} \Rightarrow~\vect{y}^*=\vect{0}$. At an equilibrium $\vect z^* = (\vect 0, \vect y^*)$, $x_i^* = s(\delta_i(\vect y^*)) = 0$ for all~$i$. Thus, \eqref{eq:update_opinion} at $\vect y^*$ evaluates, for all $i\in\mathcal V$, to be
    \begin{equation}
        y_i^* = \dfrac{\beta_i}{\beta_i+\lambda_i}\sum\nolimits_{j\in\mathcal{V}}w_{ij}y_j^*.
    \end{equation}
    Let us define $\beta_i/\beta_i+\lambda_i$ as the $i^{th}$ diagonal entry of the diagonal matrix $\mat\phi$. 
    We write the set of $y_i^*$ equilibria equations as $(\mat I- \mat{\phi}\mat{W})\vect{y}^* = \vect 0$, with
    % \begin{equation}\label{eq:y0_solution}
    %     (\mat I- \mat{\phi}\mat{W})\vect{y}^* = \vect 0.
    % \end{equation} 
    $\vect y^* = \vect 0$ obviously a solution. To prove it is in fact the unique solution, we prove the matrix $\mat{I}-\mat{\phi}\mat{W}$ is invertible, yielding $\vect{y}^* = \left(\mat{I} - \mat{\phi}\mat{W}\right)^{-1}\vect 0$. Assumption~\ref{a:0} implies that $\beta_i/(\beta_i+\lambda_i) < 1$, and thus $\mat{\phi W}$ is a nonnegative matrix with all row sums strictly less than 1. Standard nonnegative matrix theory yields that $\rho({\mat \phi W}) < 1$ (e.g. by application of \cite[Corollary 8.1.29]{horn2012matrixbook}). Thus, $\mat I - \mat{\phi W}$ cannot have a zero eigenvalue and is in fact invertible.
    % Let $\rho(\cdot)$ and $\Vert \cdot \Vert_{\infty}$ denote the spectral radius and infinity norm of a square matrix, respectively. Observe that 
    % \begin{equation}
    %     \rho(\mat{\phi W}) \leq \Vert \mat{\phi W} \Vert_\infty \leq \max_i \frac{\beta_i}{\beta_i+\lambda_i}.
    % \end{equation}
    % The right inequality is obtained by noting that $\Vert \mat W\Vert_\infty = 1$, and $\norm{\mat{\phi W}}_\infty \leq \norm{\mat{\phi}}_\infty\norm{\mat{W}}_\infty = \max_i \beta_i/(\beta_i+\lambda_i)$. Assumption~\ref{a:0} yields $\beta_i/(\beta_i+\lambda_i) < 1$, and it follows that $\rho(\mat{\phi W})<1$. Thus, $\mat I - \mat{\phi W}$ cannot have a zero eigenvalue and is in fact invertible.  
    
    Next, we show $\vect{x}^*= \vect{1} \Rightarrow~\vect{y}^*=\vect{1}$. The hypothesis that $\vect x^* = \vect 1$ implies that $s_i(\delta_i(\vect y^*))=1$ must hold for all $i \in \mathcal{V}$. 
    Hence, at an equilibrium $\vect{y}^*$, there must hold
    \begin{equation}
        y_i^* = \dfrac{\beta_i}{\beta_i+\lambda_i}\sum\nolimits_{j\in\mathcal{V}}w_{ij}y_j^* + \dfrac{\lambda_i}{\beta_i+\lambda_i}.
    \end{equation}
    The set of equilibrium equations can thus be expressed as $\vect{y}^* = \mat{\phi}\mat{W}\vect{y}^* + (\mat{I} - \mat{\phi})\vect 1$. 
    % \begin{equation}\label{eq:y1_solution}
    %     \vect{y}^* = \mat{\phi}\mat{W}\vect{y}^* + (\mat{I} - \mat{\phi})\vect 1.
    % \end{equation}
    The fact that $\mat W\vect 1 = \vect 1$ implies that $\vect y^* = \vect 1$ is a solution. We proved above that $\mat I - \mat{\phi W}$ is invertible, allowing us to rewrite $\vect y^* = (\mat I - \mat{\phi W})^{-1}(\mat{I} - \mat{\phi})\vect 1$. Clearly, the solution $\vect y^*$ is unique, and must be $\vect y^* = \vect 1$.
\end{proof}

% \begin{proposition}\label{prop:all_defect}
% Under Assumption~\ref{a:0}, the all-defection complete consensus state $\vect{z}=(\vect{0},\vect{0})$ is always an equilibrium of the coevolutionary model. Moreover, it is the unique equilibrium of the coevolutionary model if there holds
% \begin{align}\label{eq:all_D_unique_old}
%     \frac{\beta_i\lambda_i}{\beta_i+\lambda_i} \leq 2\alpha_i\left(1-\frac{r}{n}\right),
% \end{align}
% for all $i\in\mathcal V$.
% \end{proposition}
% \begin{proof}
%     To prove the first statement, we need to show that at an equilibrium $\vect z^* = (\vect x^*,\vect y^*)$, it is always true that i) $\vect x^* = \vect 0 \Rightarrow \vect y^* = \vect 0$, and ii) $\vect y^* = \vect 0 \Rightarrow \vect x^* = \vect 0$. The first part is immediate from Proposition~\ref{prop:consensus}. For the second part, observe that under the assumption $\vect y^* = \vect 0$, \eqref{eq:discriminant} evaluates to be
%     \begin{equation}
%         \delta_i(\vect y^*) = \alpha_i\bigg(\frac{r}{n}-1\bigg) - \frac{\beta_i\lambda_i}{2(\beta_i+\lambda_i)} \leq 0
%     \end{equation}  
%     because $r<n$ (recall $r$ is the public good multiplier). In other words, $x_i^* = 0$ for all $i$.

%     To prove the second statement, notice that if
%     \begin{equation}\label{eq:01}
%         \dfrac{-\alpha_i(1-\alpha_i)\left(\dfrac{r}{n}-1\right)}{\beta_i\lambda_i}+\frac{1}{2} \leq 1,
%     \end{equation}
%     then \eqref{eq:delta_neg_cond} will always hold. This would imply that $\delta_i(\vect y^*) \leq 0$ holds independently of $\vect y^*$, making $\vect z^*=(\vect 0,\vect 0)$ the unique equilibrium. Rearranging \eqref{eq:01} yields \eqref{eq:all_D_unique}. The condition in \eqref{eq:all_D_unique} shows the trade-off between self-consistency and opinion dynamics (left) and the public goods game (right).
% \end{proof}

% \begin{proposition}\label{prop:all_coop}
% Under Assumption~\ref{a:0}, the all-cooperation complete consensus state $\vect{z}=(\vect{1},\vect{1})$ is an equilibrium of the coevolutionary model if there holds
% \begin{align}\label{eq:complete_consensus_cond}
%     \frac{\beta_i\lambda_i}{\beta_i+\lambda_i} > 2\alpha_i\left(1-\frac{r}{n}\right),
% \end{align}
% for all $i\in\mathcal V$.
% \end{proposition}
% \begin{proof} 
% Similarly to Proposition~\ref{prop:all_defect}, we need to show that at an equilibrium $\vect z^* = (\vect x^*,\vect y^*)$, that i) $\vect x^* = \vect 1 \Rightarrow \vect y^* = \vect 1$, and ii) that under the condition in \eqref{eq:complete_consensus_cond} $\vect y^* = \vect 1 \Rightarrow \vect x^* = \vect 1$. The first part is the result shown in Proposition~\ref{prop:consensus}. For the second part, observe that under the assumption $\vect y^* = \vect 1$, \eqref{eq:discriminant} evaluates to be
% \begin{equation}\label{eq:02}
%         \delta_i(\vect y^*) = \alpha_i\bigg(\frac{r}{n}-1\bigg) + \frac{\beta_i\lambda_i}{2(\beta_i+\lambda_i)}.
% \end{equation} 
% For a player to cooperate, $\delta_i(\vect y^*)>0$. Setting the condition in \eqref{eq:02} to be greater than zero and rearranging yields the inequality in \eqref{eq:complete_consensus_cond}.
% \end{proof}

\begin{theorem}\label{theorem:equilibria}
    Under Assumption~\ref{a:0}, the all-defection consensus state $(\vect{0},\vect{0})$ is always an equilibrium of the coevolutionary model. Moreover, it is the unique equilibrium of the coevolutionary model if there holds
    \begin{align}\label{eq:all_D_unique}
        \frac{\beta_i\lambda_i}{\beta_i+\lambda_i} \leq 2\alpha_i\left(1-\frac{r}{n}\right),
    \end{align}
    for all $i\in\mathcal V$. If instead
       \begin{align}\label{eq:all_coop}
        \frac{\beta_i\lambda_i}{\beta_i+\lambda_i} > 2\alpha_i\left(1-\frac{r}{n}\right),
    \end{align}
    %\eqref{eq:all_D_unique} does not hold 
    for all $i\in\mathcal V$, then the all-cooperation consensus state $(\vect{1},\vect{1})$ is an equilibrium.
\end{theorem}
\begin{proof}
    To prove the first statement, we need to show that at an equilibrium $\vect z^* = (\vect x^*,\vect y^*)$, it is always true that i) $\vect x^* = \vect 0 \Rightarrow \vect y^* = \vect 0$, and ii) $\vect y^* = \vect 0 \Rightarrow \vect x^* = \vect 0$. The first part is immediate from Proposition~\ref{prop:consensus}. For the second part, observe that when $\vect y^* = \vect 0$, \eqref{eq:discriminant} evaluates to be
    \begin{equation}
        \delta_i(\vect y^*) = \alpha_i\bigg(\frac{r}{n}-1\bigg) - \frac{\beta_i\lambda_i}{2(\beta_i+\lambda_i)} \leq 0
    \end{equation}  
    because $r<n$ (recall $r$ is the public good multiplier). In other words, $x_i^* = 0$ for all $i$, hence $\vect z^* = (\vect 0,\vect 0)$ is always an equilibrium of the coevolutionary model.

    To prove the second statement, notice that if
    \begin{equation}\label{eq:01}
        \dfrac{\alpha_i(1-\alpha_i)\Big(1-\dfrac{r}{n}\Big)}{\beta_i\lambda_i}+\frac{1}{2} \geq1,
    \end{equation}
    then $\delta_i(\vect y^*) \leq 0$ holds independent of the particular value of $\vect y^*$, making $\vect z^*=(\vect 0,\vect 0)$ the unique equilibrium. Rearranging \eqref{eq:01} yields \eqref{eq:all_D_unique}. 

    To prove the all-cooperation consensus state is an equilibrium we must show that at an equilibrium $\vect z^* = (\vect x^*,\vect y^*)$, i) $\vect x^* = \vect 1 \Rightarrow \vect y^* = \vect 1$, and ii) $\vect y^* = \vect 1 \Rightarrow \vect x^* = \vect 1$. The first part is obtained from Proposition~\ref{prop:consensus}. For the second part, observe that under the assumption $\vect y^* = \vect 1$, \eqref{eq:discriminant} evaluates to be
    \begin{equation}\label{eq:02}
        \delta_i(\vect y^*) = \alpha_i\bigg(\frac{r}{n}-1\bigg) + \frac{\beta_i\lambda_i}{2(\beta_i+\lambda_i)}.
    \end{equation} 
    For a player to cooperate, $\delta_i(\vect y^*)>0$. Setting the condition in \eqref{eq:02} to be greater than zero and rearranging yields \eqref{eq:all_coop}. Therefore, if \eqref{eq:all_coop} holds for all $i \in \mathcal{V}$ then the all-cooperation consensus state is an equilibrium. %The condition in \eqref{eq:all_D_unique} shows the trade-off between self-consistency and opinion dynamics (left) and the public goods game (right).
\end{proof}

Interestingly, the conditions in \eqref{eq:all_D_unique} and \eqref{eq:all_coop} depict a trade-off between the role of self-consistency and opinion dynamics (left-hand side) and of the PGG (right-hand side). Stable cooperation seems possible when self-consistency and opinion dynamics can compensate for the fact that defection is inherently more advantageous for the single individual in the PGG. Note that when~\eqref{eq:all_D_unique} does not hold, there may be other equilibria besides the all-defection and all-cooperation consensus states, with a mixture of defectors and cooperators.

% {\color{magenta}
% \begin{remark}
%     We note that when~\eqref{eq:all_D_unique} does not hold, there are many possible equilibria; the all-defection and all-cooperation consensus states, as well as an equilibrium with clusters of defectors and cooperators.
% \end{remark}
% }

% \begin{remark}
%     We note that if only some but not all $i$ satisfy the condition for $\vect{z}^*=(\vect{1},\vect{1})$ then there are no guarantees on the uniqueness of $\vect{z}^*=(\vect{0},\vect{0})$. However, given the consensus $\vect{z}^*=(\vect{1},\vect{1})$ is not an equilibrium other non-consensus equilibria may be present.
% \end{remark}


% \begin{remark}
%     We note that the inequality in \eqref{eq:complete_consensus_cond} does not hold for all $\alpha_i,\beta_i,\lambda_i,r$ and $n$ values. Therefore unlike the equilibrium $\vect{z}^*=(\vect{0},\vect{0})$ in Proposition \ref{prop:all_defect}, the all-cooperation complete consensus state $\vect{z}^*=(\vect{1},\vect{1})$ may not always be an equilibrium.
% \end{remark}
% 

\subsection{Convergence}
% We now prove the global stability of $z^*=(0,0)$. 

Theorem~\ref{theorem:equilibria} identifies conditions on the individual-level parameters that imply that $\vect z^*=(\vect 0,\vect 0)$ is the unique equilibrium. Now, we prove that the same conditions yield convergence to said equilibrium for all initial conditions.
\begin{theorem}\label{theorem:convergence}
    Consider the coevolutionary model under Assumptions~\ref{assumption:revision_sequence} and \ref{a:0}. Suppose that $\mat W$ is symmetric and irreducible, i.e., $\mathcal G$ is undirected and connected. Suppose further that \eqref{eq:all_D_unique} holds for all $i\in\mathcal V$.
    % there holds
    % \begin{align}\label{eq:all_D_unique_con}
    %     \frac{\beta_i\lambda_i}{\beta_i+\lambda_i} \leq 2\alpha_i\left(1-\frac{r}{n}\right),
    % \end{align}
    % for all $i\in\mathcal V$.
    Then, for all $(\vect x(0),\vect y(0))$, there holds $\vect x(t)=\vect 0$ for all $t\geq T$ and $\lim_{t\to\infty}\vect y(t)=\vect{0}$, where $T$ is given in Assumption~\ref{assumption:revision_sequence}.%holds $\vect x(t) \to \vect 0$ and $\vect y(t) \to \vect 0$ as $t\to\infty$.
\end{theorem}
\begin{proof}
In this proof, we first establish convergence of $\vect x(t)$ to $\vect 0$ in finite time. Then, we show that once $\vect x(t)$ has converged, $\vect y(t)$ will converge asymptotically to $\vect 0$.

Observe that if \eqref{eq:all_D_unique} holds for all agents, then $\delta_i(\vect y(t))\leq0$, for all $i \in \mathcal{V}$ and $t\geq 0$, irrespectively of $\vect{y}(t)$. In other words, if player~$i$ is active at timestep $t$, their best-response will involve defection for the PGG (i.e., $x_i(t+1)=0$).
% will revise their strategy they will opt to defect at the next timestep. 
Under Assumption \ref{assumption:revision_sequence}, after $T$ timesteps, there holds $\vect{x}(t) = \vect 0, \forall t \geq T$. In the following, without loss of generality, we focus on analysis of $\vect y(t)$ for $t\geq T$. 

% To complete the all-defection complete consensus, we now prove $\lim_{t \to \infty} \vect y (t) = \vect 0$. 
Given that the action vector is at the all-defection action consensus state, $\vect x(t) = \vect 0$, for all $t\geq T$, the payoff that an arbitrary but fixed player~$i$ receives for having action $0$ and opinion $s\in [0,1]$ for any time $t\geq T$ simplifies to
\begin{equation}\label{eq:payoff_all_defect}
    \pi_i((0,s),\vect{z}_{-i}) = -\dfrac{1}{2}\bigg[ \beta_i \sum_{j\in \mathcal{V}} w_{ij}(s-y_j)^2 + \lambda_i s^2 \bigg].
\end{equation}
Now, consider the following potential function:
\begin{align}\label{eq:Phi_sum}
    \Phi(\vect y) = -\dfrac{1}{2} \sum_{i\in \mathcal{V}} \left[\sum_{j\in \mathcal{V}} \dfrac{w_{ij}}{2}(y_i-y_j)^2 + \dfrac{\lambda_i}{\beta_i}y_i^2\right].
\end{align}
Our hypothesis that $\mat W$ symmetric, i.e. $w_{ij}=w_{ji}$ for all $i,j \in \mathcal{V}$, means that $\Phi$ can be expressed as
\begin{align}\label{eq:Phi_sum_individual}
    &\Phi(y_i,\vect y_{-i}) = -\frac{1}{2}\bigg[\sum_{j\in \mathcal{N}_i} w_{ij}(y_i-y_j)^2 \bigg]  -\frac{1}{2}\frac{\lambda_i}{\beta_i}y_i^2 \notag \\
    &-\frac{1}{4}\sum_{k \in \mathcal{V} \setminus \{i\}}\sum_{\ell \in \mathcal{V} \setminus \{i\}} w_{k\ell}
    (y_k - y_\ell)^2 + \frac{1}{2}\bigg[ \sum_{k\in \mathcal{V}\setminus\{i\}} \frac{\lambda_k}{\beta_k}y_k^2 \bigg].
\end{align}
Notice here that we have simply rewritten $\Phi$ to highlight a fixed but arbitrary player $i \in \mathcal{V}$.

By exploiting \eqref{eq:payoff_all_defect} and \eqref{eq:Phi_sum_individual}, it can be verified that given two opinions, $s$ and $s'$ for player $i$, there holds
\begin{equation}\label{eq:payoff_potential}
    \pi(s,\vect{y}_{-i}) - \pi(s',\vect{y}_{-i}) = \beta_i(\Phi(s, \vect{y}_{-i}) - \Phi(s', \vect{y}_{-i})).
\end{equation}
Verify that the quadratic form of \eqref{eq:Phi_sum} is $\Phi(\vect y) = -\frac{1}{2} [ \vect y^\top(\mat I - \mat W + \mat{\Lambda})\vect y]$,
% \begin{equation}\label{eq:potential_matrix}
%     \Phi(\vect y) = -\frac{1}{2} \left[ \vect y^\top(\mat I - \mat W + \mat{\Lambda})\vect y \right],
% \end{equation}
where $\mat{\Lambda} = \diag(\lambda_1/\beta_1, \cdots,\lambda_n/\beta_n)$ is a positive diagonal matrix. Note that $\mat I - \mat W + \mat \Lambda$ is a symmetric nonsingular $M$-matrix~\cite[Theorem~4.31]{qu2009cooperative_book}, and thus positive definite, which means $\Phi$ has a unique maximum equal to $0$, precisely at the maximiser $\vect y = \vect 0$.

% Since $\mat W$ is symmetric, irreducible, and row-stochastic, then due to Perron–Frobenius Theorem, the eigenvalues of the matrix $\mat I - \mat W$ are real and greater than or equal to zero. That is, $\mat I - \mat W$ is positive semidefinite. The matrix $\mat{\Lambda}$ is positive definite, as it is a positive diagonal matrix. Therefore, the matrix $\mat I - \mat W + \mat{\Lambda}$ is positive definite, which means $\Phi$ has a unique maximum at $0$, precisely at the maximiser $\vect y = \vect 0$.
% . Therefore, for any $\vect y \neq 0$ we have $\Phi(\vect y)<0$. Hence, $\vect y = 0$ gives $\Phi(\vect y)=0$ and $\Phi^*=0$ must be the unique maximum to the potential function.

% Recall from Assumption \ref{assumption:revision_sequence}, there exists some timestep $T$ where all players have activated and revised their strategy at least once. 
Let player~$i$ be the active player at time $t\geq T$. From Definition~\ref{def:BR}, it is clear that the best-response dynamics implies that $\pi_i((0,y_i(t+1),\vect z_{-i}(t)) \geq \pi_i((0,y_i(t),\vect z_{-i}(t))$, with a strict inequality if $y_i(t+1) \neq y_i(t)$. This implies, from \eqref{eq:payoff_potential}, that $\Phi(\vect{y}(t+1))-\Phi(\vect{y}(t))\geq 0$, with a strict inequality if $y_i(t+1) \neq y_i(t)$. In other words, every timestep in which the active agent revises their opinion to a different value results in an increase in the potential $\Phi$. Since $\Phi(\vect y(t))$ is a monotonically non-decreasing sequence, and is bounded from above by $0$, the monotone convergence theorem~\cite{bartle1964elements} leads to the conclusion that $\lim_{t\to\infty} \Phi(\vect y(t)) = \Phi^*(\vect y^*)$. Our final step is to prove that $\Phi^* = 0$, or that $\vect y^* = \vect 0$. To see this, observe that in the limit $t\to\infty$, $\Phi(t+T) = \Phi(t)$, which is equivalent to $\vect y(t+T) = \vect y(t)$. From Assumption~\ref{assumption:revision_sequence}, this means that every agent has activated and has not changed their opinion value. In other words, $\vect y(t)$ has reached an equilibrium $\vect y^*$. We concluded earlier that $\vect x(t) = \vect 0$ for all $t\geq T$, and combined with Proposition~\ref{prop:consensus}, it follows that $\vect y^* = \vect 0$, which completes the proof.
\end{proof}

% \begin{remark}
%     We note that the coevolutionary model does not constitute a simple potential game, unless the network is undirected ($\mat W = \mat W^\top$), and all parameters are homogeneous, $\beta_i = \beta_j$ and $\lambda_i = \lambda_j$ and $\gamma_i = \gamma_j$ for all $i,j$.
% \end{remark}

\section{Conclusion}\label{sec:conclusion}

This paper introduced a mathematical model that examined the coevolution of opinions and actions in social dilemmas. We established explicit expressions for a player's best-response for the coevolutionary game and analysed the model's equilibria, focusing on the all-defection and all-cooperation consensus states. A sufficient condition for global stability of the all-defection consensus equilibrium was identified. These promising results suggest several avenues of future research. First, we aim to establish an explicit region of attraction for the all-cooperation consensus equilibrium. Second, extending from Theorem~\ref{theorem:equilibria}, we are interested in identifying conditions for the existence of equilibria with a mixture of cooperators and defectors. We also aim to show that the dynamics will always converge to an equilibrium, as extensive simulations support this conjecture.
% The coevolutionary game lies at the interface between the opinion dynamics and evolutionary game literature; each separate dynamics inherits the fundamental features of its separate grounding frameworks. 
% The player dynamics for the coevolutionary model are computed by a best-response. Through mathematical analysis, we established that the coevolutionary model experiences stable equilibria under specific conditions and has a convergence result for the all-defection consensus. Through numerical simulations, we also introduced a mixed equilibrium, which saw most players cooperating while some committed to free-riding. 

% The promising results of this paper suggest several avenues of future research. First, we aim to prove that the dynamics will always converge to an equilibrium, as extensive simulations support this conjecture. Second, extending from Theorem~\ref{theorem:equilibria}, we are interested in identifying conditions for the existence of equilibria with a mixture of cooperators and defectors (as seen in Fig.~\ref{fig:no_consensus}). Third, we aim to identify an explicit region of attraction for the all-cooperation consensus equilibrium (when it exists).
% ,  following from Proposition~\ref{eq:all_D_unique} where conditions for the all defection and all cooperation equilibria were established, defining a parameter setting when the mixed equilibrium outcome happens, as seen in Fig.~\ref{fig:no_consensus}, would be a worthy extension. Second, establishing a global convergence result for heterogenous parameters and shows the results explored in Section~\ref{sec:simulations}. 
%Finally, we hope to explore scenarios that lead to  misalignment of opinions and actions, e.g. where players support cooperation but they defect.

% mention of/tie to collective action/social dilemmas to link back to intro

\appendix

\textit{Proof of Proposition~\ref{prop:BR}:} Note that as $x_i\in\{0,1\}$, we can determine the maximisers of $\pi_i$ by evaluating the maximum of the function when $x_i=0$ and $x_i=1$, and obtain the global maximizer (i.e., the best-response) by comparing the two values obtained. We first consider the case where $x_i=1$, and differentiate $\pi_i$ with respect to $y_i$ to obtain
\begin{align}
    \pi'_i((1,y_i),\vect{z}_{-i}) = &-\beta_i(1-\gamma_i)\bigg(y_i\sum_{j\in\mathcal{V}}w_{ij} -\sum_{j\in\mathcal{V}}w_{ij}y_j\bigg) \notag\\ 
    &\qquad -\beta_i\gamma_i(y_i-u_i)-\lambda_i(y_i-1).
\end{align}
Since $\pi''_i((1,y_i),\vect z_{-i}) = -\beta_i-\lambda_i<0$, it follows that $\pi_i((1,y_i),\vect z_{-i})$ is strictly concave and has a unique maximum, allowing us to solve $\pi'_i(1,y_i) = 0$ in order to find the maximiser $y_i^1$. The same process can be followed for $x_i=0$, and we can solve $\pi'_i(0,y_i) = 0$ to find the maximiser $y_i^0$. Define $h=\beta_i(1-\gamma_i)\sum_{j\in\mathcal{V}}w_{ij}y_j+\beta_i\gamma_i u_i$ and the two maximisers can be written as
\begin{align}
    y_i^{1}=\dfrac{\beta_ih+\lambda_i}{\beta_i+\lambda_i}\;\;\text{and}\;\;
    y_i^{0}=\dfrac{\beta_ih}{\beta_i+\lambda_i}.
\end{align}
Evidently, the payoff to cooperate and defect is maximised by adopting the opinion $y_i^1$ and $y_i^0$ respectively. It follows that the best-response corresponds to the greater of $\pi_i((1,y_i^{1}),\vect{z}_{-i})$ versus $\pi_i((0,y_i^{0}),\vect{z}_{-i})$. Verify that the discriminant $\delta_i(z_i,\vect{z}_{-i})$ in \eqref{eq:discriminant} is precisely $\pi_i((1,y_i^{1}),\vect{z}_{-i})-\pi_i((0,y_i^{0}),\vect{z}_{-i})$.
% By computing
% \begin{align}\label{eq:final_difference}
%     &\pi_i((+1,y_i^{+1}),\vect{z}_{-i})-\pi_i((0,y_i^{0}),\vect{z}_{-i}) = \notag \\ 
%     &\alpha_i\left( \dfrac{r}{n}-1 \right) + \dfrac{\lambda_i h + \frac{\beta_i\lambda_i}{2}}{\beta_i+\lambda_i},
% \end{align}
% it follows that the discriminant $\delta_i(z_i,\vect{z}_{-i})$ in \eqref{eq:discriminant} is precisely $\pi_i((1,y_i^{1}),\vect{z}_{-i})-\pi_i((0,y_i^{0}),\vect{z}_{-i})$. 
Thus, $\delta_i(z_i,\vect{z}_{-i}) > 0$ and $\delta_i(z_i,\vect{z}_{-i}) < 0$ indicate that $(1,y_i^{1})$ and $(0,y_i^{0})$ are the best-response strategies, respectively. If $\delta_i(z_i,\vect{z}_{-i}) = 0$, both $(1,y_i^{1})$ and $(0,y_i^{0})$ are best-responses. \hfill $\qed$

\bibliographystyle{IEEEtran}
\bibliography{ref}

\end{document}
